% Propietario conservado: /Users/ruben/Documents/New project/output/EXTREMA_MEDIA_RAZON_RECIPROCIDAD_ES_EN_20260918/ARTICULO/ES/source/libro/05_levantamiento.tex
\chapter{Unitary lift of refinement and transport of readers}
\label{x_reciprocidad:lib05:levantamiento}

The residue bijection of \cref{x_reciprocidad:lib03:teo-residuo} and the
bilateral balance of \cref{x_reciprocidad:lib04:teo-balance} admit an
explicit composition. An admissible arithmetic triple, containing
a pair and its carry, determines a refined pair. The
linearization of this bijection on orthonormal bases provides
a unitary transport; two transports of this type can act
as the two realizations of the bilateral balance.

The realization is established after the objects generated by
APP--TRIT--TPK. Its combinatorial domain comprises the pairs and
carries satisfying the conditions written below. The norm
preserved by its linearization is the sum of squared amplitudes
over those labels; the arithmetic conservation of
\(L_c(x,y)=(Bx-c,By+c)\) remains the sum of the two
coordinates after rescaling. The composition relates both
conservations through a specific map.

\section{Bijection of admissible arithmetic states}
\label{x_reciprocidad:lib05:biyeccion}

Let \(B\geq2\) and \(M\geq1\) be integers. We define
\begin{align}
 \mathcal D_{B,M}
 &=\bigl\{(x,y,c)\in\mathbb Z_{\geq0}^2
                 \times\{0,\ldots,B-1\}:x+y=M,\ c\leq Bx\bigr\},
 \label{x_reciprocidad:lib05:dominio}\\
 \mathcal E_{BM}
 &=\bigl\{(X,Y)\in\mathbb Z_{\geq0}^2:X+Y=BM\bigr\}.
 \label{x_reciprocidad:lib05:codominio}
\end{align}
The condition \(c\leq Bx\) maintains the nonnegativity of the
first output coordinate. For \(x=0\) it only admits
\(c=0\); for \(x\geq1\) it admits the entire canonical alphabet.

\begin{theorem}[Exact correspondence of labels]
\label{x_reciprocidad:lib05:teo-biyeccion}
The map
\begin{equation}
 \ell_{B,M}:\mathcal D_{B,M}\longrightarrow\mathcal E_{BM},
 \qquad (x,y,c)\longmapsto(Bx-c,By+c)
 \label{x_reciprocidad:lib05:mapa-etiquetas}
\end{equation}
is bijective. Both sets have cardinality \(BM+1\). Its inverse
is given by
\begin{equation}
 c=Y\bmod B,\qquad y=\frac{Y-c}{B},\qquad x=\frac{X+c}{B}.
 \label{x_reciprocidad:lib05:inversa-etiquetas}
\end{equation}
\end{theorem}

\begin{proof}
The domain condition and the identity
\(Bx-c+By+c=BM\) ensure that the image belongs to
\(\mathcal E_{BM}\). Given a pair in that set, the residue
\(c\) is the unique alphabet element satisfying
\(Y\equiv c\pmod B\). Since \(X+Y\) is divisible by \(B\),
so is \(X+c\). Euclidean division of \(Y\) gives
\(y\geq0\), and \(X+c\geq0\) gives \(x\geq0\). If \(c>0\),
\(x=0\) would require \(X=-c<0\), so admissibility is
preserved. The sum of the preimages is \(M\), and
substitution verifies the original image. Uniqueness of the residue and
quotients proves the bijection.

There are \(M+1\) labels with \(c=0\) and \(M\) with each positive
carry; therefore the cardinality is
\((M+1)+(B-1)M=BM+1\), equal to that of the output.
\end{proof}

The proof makes explicit the boundary endpoint of the domain, a necessary
condition for equality of cardinalities. The new carry label
is a typed part of the input to refinement.

\section{Unitarity on label amplitudes}
\label{x_reciprocidad:lib05:unitariedad}

Take the real or complex Hilbert spaces with counting measure
\begin{equation}
 \mathscr H_D=\ell^2(\mathcal D_{B,M}),\qquad
 \mathscr H_E=\ell^2(\mathcal E_{BM}),
 \label{x_reciprocidad:lib05:hilbert}
\end{equation}
and their labeled orthonormal bases. Define linearly
\begin{equation}
 J_{B,M}|x,y,c\rangle=|Bx-c,By+c\rangle.
 \label{x_reciprocidad:lib05:j}
\end{equation}
We shall write \(J\) when the parameters are fixed.

\begin{proposition}[Unitary lift]
\label{x_reciprocidad:lib05:prop-unitaria}
The operator \(J:\mathscr H_D\to\mathscr H_E\) is unitary:
\(J^*J=I_D\), \(JJ^*=I_E\). The adjoint recovers each label
through \eqref{x_reciprocidad:lib05:inversa-etiquetas}.
\end{proposition}

\begin{proof}
\cref{x_reciprocidad:lib05:teo-biyeccion} ensures that \(J\) maps
two complete orthonormal bases bijectively. For
\(\psi=\sum_{d\in\mathcal D}\psi_d|d\rangle\),
\[
 \|J\psi\|^2
 =\sum_{e\in\mathcal E}|\psi_{\ell_{B,M}^{-1}(e)}|^2
 =\sum_{d\in\mathcal D}|\psi_d|^2=\|\psi\|^2.
\]
Surjectivity proves the two unitary identities and the
inverse bijection describes the adjoint on each basis vector.
\end{proof}

The label \(|73,27,1\rangle\) has norm one; the integers
\(73,27,1\) describe its arithmetic values, not its amplitudes.
The construction represents reversible processing of labels
or amplitudes. A specific physical realization additionally preserves
the protocol preparing and transporting those amplitudes.

\section{Diagonal proportion and carry readers}
\label{x_reciprocidad:lib05:diagonales}

For a real function \(h\) on a finite set, denote
by \(M_h\) its diagonal operator: each basis vector is multiplied
by the value of \(h\) at its label. The bijection satisfies
\begin{equation}
 J^*M_hJ=M_{h\circ\ell_{B,M}}.
 \label{x_reciprocidad:lib05:lector-diagonal}
\end{equation}
This is proved by acting on all basis vectors.

Define on the output
\begin{equation}
 f_E(X,Y)=\frac{X}{BM},\quad
 g_E(X,Y)=\frac{Y}{BM},\quad
 \gamma_E(X,Y)=Y\bmod B,
 \label{x_reciprocidad:lib05:lectores-salida}
\end{equation}
and on the input \(f_D=x/M\), \(g_D=y/M\), \(\gamma_D=c\).
The identity \eqref{x_reciprocidad:lib05:lector-diagonal} provides
\begin{align}
 J^*M_{f_E}J&=M_{\,x/M-c/(BM)},\nonumber\\
 J^*M_{g_E}J&=M_{\,y/M+c/(BM)},\nonumber\\
 J^*M_{\gamma_E}J&=M_c.
 \label{x_reciprocidad:lib05:transporte-fracciones}
\end{align}
The first two readers transport the unit transfer.
Since \(f_E+g_E=1\), their sum is \(I_E\) and their compression is
\(I_D\). The third reader exactly preserves the carry as
an integer observable, with its residual representation in the output.

\section{Complete routes at arbitrary depth}
\label{x_reciprocidad:lib05:rutas}

Let \(\mathcal P_n(B,M)\) be the set of routes of length \(n\)
starting from all nonnegative pairs with sum \(M\) and, at
each stage, satisfying the conditions of
\eqref{x_reciprocidad:lib05:dominio} for the corresponding sum. A route preserves
its initial pair and its \(n\) admissible carries. Let
\(\Phi_n\) be the map assigning each route its final pair.

\begin{theorem}[Recovery of routes at every finite depth]
\label{x_reciprocidad:lib05:teo-rutas}
For every \(n\geq0\),
\begin{equation}
 \Phi_n:\mathcal P_n(B,M)\longrightarrow\mathcal E_{B^nM}
 \quad\hbox{is bijective},\qquad
 |\mathcal P_n(B,M)|=B^nM+1.
 \label{x_reciprocidad:lib05:biyeccion-rutas}
\end{equation}
Its linearization
\begin{equation}
 J_n:\ell^2(\mathcal P_n(B,M))\longrightarrow
          \ell^2(\mathcal E_{B^nM}),\qquad
 J_n|\pi\rangle=|\Phi_n(\pi)\rangle
 \label{x_reciprocidad:lib05:jn}
\end{equation}
is unitary.
\end{theorem}

\begin{proof}
From a pair with sum \(B^nM\), the residue of
\eqref{x_reciprocidad:lib05:inversa-etiquetas} recovers the last carry and a
unique preceding pair with sum \(B^{n-1}M\). Repeating \(n\) times
produces a unique route whose input sums to \(M\). The direct
composition recovers the original endpoint, proving bijection and
inverse. The count is also obtained from the inputs: the
\(M\) pairs with \(x_0\geq1\) admit \(B^n\) words each,
whereas \((0,M)\) admits only the all-zero
word. The total is \(MB^n+1\). The bijective transformation of
orthonormal bases proves the unitarity of \(J_n\).
\end{proof}

With the depth known, the readers on the final pair are
\begin{equation}
 c_j=\left\lfloor\frac{Y}{B^{n-j}}\right\rfloor\bmod B,
 \quad y_0=\left\lfloor\frac{Y}{B^n}\right\rfloor,
 \quad x_0=M-y_0,
 \qquad 1\leq j\leq n.
 \label{x_reciprocidad:lib05:lectores-ruta}
\end{equation}
These functions are diagonal and recover the complete memory of
the finite operations. Compatibility with prolongation is
\begin{equation}
 \Phi_{n+1}(\pi,c)=L_c(\Phi_n(\pi)).
 \label{x_reciprocidad:lib05:compatibilidad-prolongacion}
\end{equation}
On passing from \(n\) to \(n+1\), the new admissible carry
label is incorporated. The spaces increase in dimension and the
boundary label admits only \(c=0\); therefore the domain of the
next step preserves admissibility, instead of replacing it
with an unrestricted product with \(B\) digits.

\section{Fixed input and exact image}
\label{x_reciprocidad:lib05:entrada-fija}

If an initial pair \((x_0,y_0)\) is fixed, with \(x_0\geq1\),
there are \(B^n\) routes and their image is exactly
\begin{equation}
 \bigl\{(B^nx_0-C,B^ny_0+C):0\leq C\leq B^n-1\bigr\}.
 \label{x_reciprocidad:lib05:imagen-fija}
\end{equation}
\cref{x_reciprocidad:lib03:cor-palabra} proves that the accumulated \(B^n\)
are distinct and recover the respective words. The
linearization is unitary onto the subspace spanned by this
image. If \(x_0=0\), the image consists of a single label.

The cylinders of \cref{x_reciprocidad:lib03:teo-cilindros} correspond to a
fixed input. The complete domain of
\eqref{x_reciprocidad:lib05:biyeccion-rutas} brings together all pairs with sum \(M\).
Both constructions have explicit inverses, but their images
and cardinalities differ. Arbitrary depth expresses the
validity of the proof for every \(n\); the limit of the reading and
its boundaries remain as established in
\cref{x_reciprocidad:lib03:limite}.

\section{Bilateral coupling of refinement}
\label{x_reciprocidad:lib05:bilateral}

Let \(R\) be a previously declared permutation of
\(\mathcal D_{B,M}\), unitarily linearized on
\(\mathscr H_D\). Define the two transports
\begin{equation}
 \mathcal B=J,\qquad \mathcal C=JR:
            \mathscr H_D\longrightarrow\mathscr H_E.
 \label{x_reciprocidad:lib05:dos-transportes}
\end{equation}
The permutation is part of the specification of the procedure
and is fixed before evaluating an output reading.
For \(a>0\), \(p=a^2+a+1\), \(N=(2a+1)p\),
\begin{equation}
 Q_-=J(aI-R),\qquad Q_+=J((a+1)I+R).
 \label{x_reciprocidad:lib05:q-refinamiento}
\end{equation}

\begin{proposition}[Distribution and recovery of refined amplitudes]
\label{x_reciprocidad:lib05:prop-bilateral}
The operators \eqref{x_reciprocidad:lib05:q-refinamiento} satisfy
\begin{equation}
 (a+1)Q_-^*Q_-+aQ_+^*Q_+=NI_D,
 \label{x_reciprocidad:lib05:balance}
\end{equation}
and the map
\begin{equation}
 W_a\psi=\frac1{\sqrt N}
    \bigl(\sqrt{a+1}\,Q_-\psi,\sqrt a\,Q_+\psi\bigr)
 \label{x_reciprocidad:lib05:w}
\end{equation}
is isometric. For the unnormalized channels
\(u=Q_-\psi\), \(v=Q_+\psi\),
\begin{equation}
 J\psi=\frac{u+v}{2a+1},\qquad
 \psi=J^*\frac{u+v}{2a+1}.
 \label{x_reciprocidad:lib05:inversa-bilateral}
\end{equation}
\end{proposition}

\begin{proof}
The two maps \(\mathcal B,\mathcal C\) are unitary.
The cross terms of the Gram matrices of \(Q_-\) and \(Q_+\) have
coefficients \(-a\) and \(a+1\), respectively. The balance
weights cancel them, while the diagonal terms sum to
\(N\), as in \cref{x_reciprocidad:lib04:teo-balance}. Normalization
proves isometry. The sum of the channels is
\((2a+1)J\psi\), and the unitarity of \(J\) gives the inverse.
\end{proof}

From the normalized channels the input is recovered directly
by \(W_a^*\), with inversion norm one; this adjoint already incorporates
\(J^*\). The alternative formula
\eqref{x_reciprocidad:lib05:inversa-bilateral} separates recovery of \(J\psi\)
by summation from application of \(J^*\), whose action on labels
is calculated through residues. Reversible arithmetic refinement
and bilateral amplitude transport are thereby composed.

\begin{proposition}[Compatibility in the refined space]
\label{x_reciprocidad:lib05:prop-compatibilidad}
An unnormalized pair \((u,v)\) belongs to the image of
\(\psi\mapsto(Q_-\psi,Q_+\psi)\) if and only if
\begin{equation}
 \bigl[(a+1)I_E+JRJ^*\bigr]u
 =\bigl[aI_E-JRJ^*\bigr]v.
 \label{x_reciprocidad:lib05:compatibilidad}
\end{equation}
\end{proposition}

\begin{proof}
The equality is equivalent to
\(JRJ^*(u+v)=av-(a+1)u\). If it holds, define
\(\psi=J^*(u+v)/(2a+1)\). On substitution into
\eqref{x_reciprocidad:lib05:q-refinamiento}, this identity gives
\(Q_-\psi=u\) and \(Q_+\psi=v\). Necessity follows from
inserting the true channels. All compositions act
on \(\mathscr H_E\), so their types are defined.
\end{proof}

\section{Law of observables on the two channels}
\label{x_reciprocidad:lib05:ley-lectores}

\begin{theorem}[Bilateral compression of a refined reader]
\label{x_reciprocidad:lib05:teo-lector}
For \(F=F^*\) on \(\mathscr H_E\), write
\(F_J=J^*FJ\). Then
\begin{equation}
 W_a^*\operatorname{diag}(F,F)W_a
 =\frac{a(a+1)F_J+R^*F_JR}{p}.
 \label{x_reciprocidad:lib05:lector-bilateral}
\end{equation}
The reader \(F_J\) is preserved exactly if and only if it commutes
with \(R\).
\end{theorem}

\begin{proof}
Expanding \((a+1)Q_-^*FQ_-+aQ_+^*FQ_+\) cancels the
mixed terms. The coefficient of \(\mathcal B^*F\mathcal B\)
is \(a(a+1)(2a+1)\); that of
\(\mathcal C^*F\mathcal C\) is \(2a+1\). Dividing by
\(N\), substituting \(\mathcal C=JR\) and using
\(F_J=J^*FJ\) proves the formula. Equality with \(F_J\)
is equivalent to \(R^*F_JR=F_J\), or to \([F_J,R]=0\) by
unitarity.
\end{proof}

In the case \(a=8\), the weights are \(72/73\) and \(1/73\).
The identity reader recovers the norm. A nonconstant reader
may change even when the norm is preserved; its variation is
calculated by \eqref{x_reciprocidad:lib05:lector-bilateral}.
The operator that exactly preserves an original reader \(A\)
in the complete image is \(W_aAW_a^*\), whose image identity
is \(W_aW_a^*\). Reading \(\operatorname{diag}(F,F)\) corresponds
to the specific compression of the theorem and does not automatically
replace that transported observable.

\section{Arithmetic example of dimension one thousand and one}
\label{x_reciprocidad:lib05:ejemplo}

Take \(B=10\), \(M=100\). The label spaces have
dimension \(1001\), and
\begin{equation}
 d_1=(73,27,1)\longmapsto e_1=(729,271).
 \label{x_reciprocidad:lib05:ejemplo-etiqueta}
\end{equation}
The diagonal readers give
\begin{equation}
 f_E(e_1)=\frac{729}{1000}
 =\frac{73}{100}-\frac1{1000},\qquad
 g_E(e_1)=\frac{271}{1000},\qquad \gamma_E(e_1)=1.
 \label{x_reciprocidad:lib05:ejemplo-lectores}
\end{equation}
Declare \(R\) to be the transposition of the admissible
labels \(d_1=(73,27,1)\) and \(d_2=(73,27,2)\), leaving
the others fixed. The second image is \(e_2=(728,272)\).
The transposition rule has been fixed on the labels, before
any comparison with a constant.

For \(a=8\) and \(\psi=|d_1\rangle\),
\begin{equation}
 Q_-\psi=8|e_1\rangle-|e_2\rangle,\qquad
 Q_+\psi=9|e_1\rangle+|e_2\rangle.
 \label{x_reciprocidad:lib05:ejemplo-canales}
\end{equation}
The squared norms are \(65\) and \(82\), and
\begin{equation}
 9\cdot65+8\cdot82=585+656=1241=17\cdot73.
 \label{x_reciprocidad:lib05:ejemplo-balance}
\end{equation}
The sum of the channels is \(17|e_1\rangle\). The inverse
\eqref{x_reciprocidad:lib05:inversa-bilateral} and the residue recover
\(|d_1\rangle\), including the carry one.

If \(F=M_{f_E}\) is evaluated diagonally in both normalized
channels, \cref{x_reciprocidad:lib05:teo-lector} gives
\begin{equation}
 \begin{split}
 \langle W_8\psi,\operatorname{diag}(F,F)W_8\psi\rangle
 &=\frac{72\cdot729+728}{73\cdot1000}\\
 &=\frac{729}{1000}-\frac1{73000}.
 \end{split}
 \label{x_reciprocidad:lib05:ejemplo-fraccion}
\end{equation}
The complementary reading is one minus this quantity. For the
diagonal carry reader we obtain
\begin{equation}
 \frac{72\cdot1+2}{73}=\frac{74}{73}.
 \label{x_reciprocidad:lib05:ejemplo-acarreo}
\end{equation}
The original carry remains exactly recoverable: the
observable \(W_8M_cW_8^*\) has on \(W_8|d_1\rangle\)
the eigenvalue one. The value \(74/73\) is the response of the
diagonal reader applied to both channels, whose relative transport
is specified by \eqref{x_reciprocidad:lib05:lector-bilateral}. The example
distinguishes conservation of information from variation of a particular
reading.

\section{Articulation with the generated register and application domains}
\label{x_reciprocidad:lib05:dominios}

The unitary \(J\) acts on \(BM+1\) labels, whereas
the realization of \(K\) has twelve coordinates. In the example,
\(1001\) counts refined pairs with sum one thousand, and twelve counts
windows of the produced register. Both dimensions retain their
specific combinatorial origin.

A particular composition between these two carriers can be formulated
through a material selector of twelve distinct routes, with its
provenance and operators. If that selector provides an
isometry
\(J_{12}:\mathbb R^{12}\to\ell^2(\mathcal D_{B,M})\), then
\(JJ_{12}\) is isometric, and the intertwining of the readers is
checked on it. This is the hypothesis of that additional
composition, not an implicit identification through the number of
coordinates. The bilateral maps already defined directly
on \(K\) in \cref{x_reciprocidad:lib04:bilateral} retain their complete
domains without using that additional selector.

The result of this section brings together an arithmetic bijection, its
unitary, intertwined diagonal readers, route recovery at
any finite depth and an explicit bilateral composition.
The route domain is that defined in
\eqref{x_reciprocidad:lib05:biyeccion-rutas}. Applications to histories
filtered by TPK retain their survival and
prolongation rules; the label correspondence here proves the
declared transport, without replacing those operators with an unrestricted
arithmetic domain.
